\documentclass[11pt,a4paper]{article}

% Packages
\usepackage[utf8]{inputenc}
\usepackage[T1]{fontenc}
\usepackage[margin=0.75in]{geometry}
\usepackage{enumitem}
\usepackage{hyperref}
\usepackage{titlesec}
\usepackage{xcolor}

% Formatting
\pagestyle{empty}
\setlength{\parindent}{0pt}
\setlist[itemize]{leftmargin=*, noitemsep, topsep=3pt}

% Section formatting
\titleformat{\section}{\Large\bfseries}{}{0em}{}[\titlerule]
\titlespacing*{\section}{0pt}{12pt}{6pt}

\titleformat{\subsection}[runin]{\bfseries}{}{0em}{}
\titlespacing*{\subsection}{0pt}{8pt}{0pt}

% Colors
\definecolor{linkcolor}{RGB}{0,102,204}
\hypersetup{
    colorlinks=true,
    urlcolor=linkcolor
}

\begin{document}

% Header
\begin{center}
    {\LARGE \textbf{ERIC MAYEUX}}\\
    \vspace{4pt}
    {\large Gestionnaire de projet TI | IT Project Manager}\\
    \vspace{4pt}
    Montréal, QC, Canada\\
    \href{mailto:mayeux.e@gmail.com}{mayeux.e@gmail.com} | 438-884-7645
\end{center}

\vspace{8pt}

\section*{Profil Professionnel}
\noindent
\textbf{Gestionnaire de projet TI} avec 10+ ans d'expérience en \textbf{pilotage de projets} dans le secteur bancaire et financier. Expert en \textbf{planification, coordination et suivi de performance} de projets TI complexes, incluant des implémentations ERP/logicielles. Maîtrise des méthodologies \textbf{Agile/Scrum} et des outils (Jira, Azure DevOps, MS Project). Certifié \textbf{PMP} (en cours), CSM, ITIL [À VÉRIFIER]. Bilinguisme français/anglais. Orientation résultats et rigueur organisationnelle.

\section*{Compétences Techniques}

\textbf{Gestion de projets TI:} Planification, coordination, calendriers de projet, suivi d'exécution, budgets, gestion des risques, plans de contingence, livraison logicielle (ERP, CRM)

\textbf{Leadership \& Communication:} Coordination multi-parties prenantes, résolution de problèmes, communication technique/non-technique, bilingue

\textbf{Gestion du changement:} Intégration des systèmes, transitions BAU, suivi des fournisseurs, normes de sécurité, gouvernance

\textbf{Outils:} Jira, Azure DevOps, MS Project, Confluence, Datadog, Splunk, tableaux de bord KPI

\textbf{Méthodologies:} Agile/Scrum, SAFe, Kanban, Waterfall, hybride, PMP

\textbf{Intelligence Artificielle:} Agents IA, Playwright MCP, Crew AI, automatisation, RPA

\textbf{Architecture \& Cloud:} AWS Cloud Practitioner, intégration de systèmes, Big Data

\vspace{6pt}

\section*{Réalisations}
\begin{itemize}
    \item \textbf{Pilotage de projets TI} à la BNC : coordination de A à Z, suivi des livrables et budgets
    \item Mise en place de \textbf{calendriers de projet fiables} et suivi des indicateurs de performance
    \item \textbf{Gestion des risques} proactive : plans de contingence, actions correctives, gestion des fournisseurs
    \item Coordination multi-sites (Thaïlande) : \textbf{intégration de systèmes} et gestion des dépendances
    \item Créé des \textbf{tableaux de bord} pour le suivi de la santé des projets et la prise de décision
    \item Déploiement d'\textbf{agents IA} pour l'automatisation des tâches de gestion TI
\end{itemize}

\section*{Expérience Professionnelle}

\subsection*{Scrum Master -- Gestion de projet | Project Manager} \hfill \textit{Novembre 2025 -- Présent}\\
\textbf{Banque Nationale du Canada (BNC)} -- Montréal, QC\\
\textbf{Pilote} les projets TI : planification, coordination, suivi des livrables, gestion des risques.
\begin{itemize}
    \item Élabore et maintient les \textbf{calendriers de projet} ; surveille l'exécution des initiatives TI
    \item Coordonne les exigences et \textbf{livrables TI} avec les parties prenantes internes et externes
    \item Réalise des \textbf{évaluations des risques} et élabore des plans de contingence
    \item Assure la \textbf{bonne intégration des systèmes} d'information dans les projets
    \item Prépare les \textbf{rapports et bilans post-implantation} ; supervise les transitions BAU
    \item Développe des \textbf{agents IA} pour l'automatisation des tâches de gestion TI
\end{itemize}

\subsection*{Intégrateur Sénior Qualité | Senior Quality Integrator} \hfill \textit{Octobre 2023 -- Novembre 2025}\\
\textbf{Banque Nationale du Canada (BNC)} -- Montréal, QC\\
\begin{itemize}
    \item \textbf{Gestion de projets TI} multi-sites : planification, livrables, gestion des risques
    \item \textbf{Intégration de systèmes} : CI/CD, pipelines de test, Jira, Datadog, Splunk
    \item Coordination de fournisseurs et prestataires externes (Thaïlande)
\end{itemize}

\subsection*{Intégrateur Qualité -- SDET | Quality Integrator -- Software Development Engineer in Test} \hfill \textit{Janvier 2022 -- Octobre 2023}\\
\textbf{Banque Nationale du Canada (BNC)} via Groupe Astek -- Montréal, QC\\
\begin{itemize}
    \item \textbf{Implémentation logicielle} : stratégie de test, métriques qualité, sécurité
    \item Stack : Selenium, AppliTools, RestAssured, JMeter, Gatling, K6
\end{itemize}

\subsection*{QA Automaticien | QA Automation Engineer} \hfill \textit{Juin 2021 -- Décembre 2021}\\
\textbf{AODocs} -- Paris, France\\
\begin{itemize}
    \item Automatisation tests non-régression ; POC Flutter ; rapports Allure
\end{itemize}

\subsection*{Automaticien CI/CD | DevOps Automation Engineer} \hfill \textit{Janvier 2020 -- Juin 2021}\\
\textbf{ENGIE DIGITAL} via Groupe Hénix -- Paris, France\\
\begin{itemize}
    \item Création CI/CD : Jenkins, Docker ; \textbf{administration Jira}, KPIs, scripting Python/Shell
\end{itemize}

\subsection*{Product Owner} \hfill \textit{Juin 2019 -- Septembre 2019}\\
\textbf{Hénix} -- Montrouge, France\\
\begin{itemize}
    \item Gestion du backlog, priorisation ; certification Jira ACP-600 (8 500 utilisateurs)
\end{itemize}

\subsection*{Développeur Logiciel | Software Developer} \hfill \textit{Mai 2017 -- Mai 2019}\\
\textbf{MEDIAPOST} via Groupe Hénix -- Paris, France\\
\begin{itemize}
    \item Développement back office, webservices SOAP/REST, \textbf{intégration de systèmes}
\end{itemize}

\subsection*{Consultant QA \& Automatisation | QA Automation Consultant} \hfill \textit{Avril 2016 -- Avril 2017}\\
\textbf{ENGIE IT} via Groupe Hénix -- Saint-Ouen, France\\
\begin{itemize}
    \item \textbf{Gestion des fournisseurs}, administration HP ALM/UFT, sécurité des environnements
\end{itemize}

\section*{Certifications \& Formations}

\textbf{PMP Certification} (en cours) \hfill \textit{2026}

Project Management Professional -- Project Management Institute (PMI)

\textbf{AWS Cloud Practitioner} \hfill \textit{2026}

Amazon Web Services (AWS)

\textbf{Certified Scrum Master (CSM)} \hfill \textit{2024}

Scrum Alliance

\textbf{Jira ACP-600 -- Project Administration in Jira Server} \hfill \textit{2019}

Atlassian Certified Professional

\textbf{Cursus Automaticien et DevOps} \hfill \textit{2019}\\
\textit{AFCEPF -- Montrouge, France}

\textbf{Cursus Architecture logiciel \& Analyste informaticien - Certifié Master II} \hfill \textit{2017}\\
\textit{AFCEPF-HENIX -- Malakoff, France}

\section*{Formation Académique}

\textbf{Licence en Gestion (Bachelor's Degree in Management)} \hfill \textit{2011}\\
École Supérieure de Gestion (E.S.G) -- Paris, France

\section*{Compétences Linguistiques}

\textbf{Français:} Langue maternelle (Native)\\
\textbf{Anglais:} Niveau Avancé (Advanced / Professional Working Proficiency)

\end{document}
